\section{Implications for Contract Design}\label{sec:implications}

If the 2026 wave raises the perceived probability of demands ($q$), Proposition~\ref{prop:price} predicts that protection will be repriced. The model also says which features of a clause matter most, so it predicts the form the repricing should take.

\begin{enumerate}[leftmargin=*,itemsep=2pt]
\item \textbf{Ban lists that are too short, not lists themselves.} Because order statistics flatten quickly, a no-trade list that leaves 15 or 20 approved teams costs the team almost nothing on exit. A trade list naming three costs it most of the value. The efficient contract caps the player's veto at a minimum approved set (for example, at least eight teams). It does not have to deny protection altogether.
\item \textbf{Schedule the clause to decay.} Converting an NMC into a modified NTC on a fixed date, and having lists shrink over time, raises $m$ in later years. That makes waiting less attractive to the player and selling more attractive to the team (Proposition~\ref{prop:delay}). Many recent contracts already do this, and more should.
\item \textbf{Stagger clause dates across the roster.} A team whose stars' protection windows open and close together has no flexibility in a crisis. Staggering keeps some assets always tradable.
\item \textbf{Price the clause explicitly.} The absence of a measurable wage discount (Section~\ref{sec:empirics}) suggests that teams have treated protection as free. A team that values the option to sell should demand cap savings for an NMC, or offer an NMC only in the years when demands are least likely.
\item \textbf{Keep requests private.} Since audience costs destroy the seller's outside option only when the approved list is short, teams have the strongest interest in confidentiality exactly when a player names few teams. A protocol committing both sides to a private negotiation window before any public statement would protect the seller's surplus without removing the player's right.
\item \textbf{Bargain over it at the league level.} Bettman's remark makes clear that the clause will be on the agenda for the next CBA. The model suggests a reform that protects the players' interest in stability while restoring competition at exit. One example: a league-wide minimum for the size of approved lists after a player-initiated request. It need not restrict clause eligibility.
\end{enumerate}

The argument also generalises beyond hockey. Any veto right granted to protect a weaker party in a long-term relationship can, once the relationship sours, become an instrument with which that party, or the third party it favours, extracts surplus. Examples include consent rights in joint ventures, anti-dilution and tag-along rights for minority investors, and alliance commitments that require partner approval before forces move. In each case the right is ex ante protective and ex post extractive. In each case the design remedy is the one proposed here: limit how far the veto can narrow the field, not whether the veto exists.
